\section*{Author Biographies}

\newcommand{\paperauthorbio}[3]{%
\par\medskip\noindent
\begin{minipage}[t]{0.22\columnwidth}
\vspace{0pt}
\includegraphics[width=\linewidth,height=1.12in,keepaspectratio]{#1}
\end{minipage}\hfill
\begin{minipage}[t]{0.74\columnwidth}
\vspace{0pt}
\textbf{#2} #3
\end{minipage}
\par\medskip
}

\paperauthorbio{figures/authors/mayank_lohani.png}{Mayank Lohani}{is currently pursuing a Bachelor of Technology degree in Computer Science and Engineering with specialization in Artificial Intelligence and Machine Learning at the School of Computer Science and Engineering (SCOPE), Vellore Institute of Technology, Chennai, India. He is expected to graduate in 2028. His areas of interest include data science, machine learning, artificial intelligence, and their applications in solving real-world problems.}

\paperauthorbio{figures/authors/sudha_a.jpg}{Sudha Anbalagan}{received the Ph.D. degree from the Department of Information Technology, Anna University, Chennai, India. She was an IT Analyst with Tata Consultancy Services Ltd. She was also a Visiting Research Fellow with the Department of Computer Science, University of California, Davis, Davis, USA, in 2015. She is currently an Associate Professor with the Centre for Smart Grid Technologies, School of Computer Science and Engineering, Vellore Institute of Technology Chennai, Chennai, India. Her research interests include energy optimization, vehicular networks, UAVs, the Internet of Things, software-defined networking, network security, and machine learning. She was a recipient of the Anna Centenary Research Fellowship during her Ph.D. program.}

\paperauthorbio{figures/authors/prabhakar_m.jpg}{M. Prabhakar}{received the B.E. degree in electrical and electronics engineering from the University of Madras, Chennai, India, in 1998, the M.E. degree in power electronics and drives from Bharathidasan University, Tiruchirappalli, India, in 2000, and the Ph.D. degree in electrical engineering from Anna University, Chennai, India, in 2012. He started his teaching career as a Lecturer in 2000. Since 2012, he has been an Associate Professor with the School of Electrical Engineering (SELECT), Vellore Institute of Technology Chennai, Chennai, India. Since 2019, he has been working as a Professor and has been associated with the Centre for Smart Grid Technologies since May 2022. He has coauthored about 45 research papers in reputed journals and conferences. His research interests include power electronics, power converters, high-gain dc-dc converters, and dc microgrids. He was a recipient of the Outstanding Teacher Award for his excellent teaching and research contributions in 2009 and the Research Award from Vellore Institute of Technology for his research contributions continuously since 2012. He is an active reviewer for reputed journals.}

\paperauthorbio{figures/authors/arul_r.jpg}{Arul R.}{received the B.E., M.E., and Ph.D. degrees in electrical and electronics engineering from Annamalai University, India, in 1995, 1996, and 2016, respectively. He is currently a Faculty Member at the School of Electrical Engineering, Vellore Institute of Technology Chennai, Chennai, India. He has 24 years of teaching and research experience, has published more than 60 international journal articles indexed in Scopus and Web of Science, and holds ten patents. His research interests include power-system optimization, soft-computing techniques, the Internet of Things, and electric vehicles.}
